% GENERATED FILE. Edit canonical lessons, then run npm run generate:books.
% canonical-source-hash: fnv1a64:b964ea1554398082
% canonical-lessons: TE-C152-jhankaram, TE-S173-letter-jha, TE-R152-jhankaram-recall

\chapter{A Word for Buzzing — Then Its First Letter}
\label{ch:te-jhankaram-jha}
\hlchaptermodality{152}

\begin{chapteropening}
\textbf{By the end of this chapter:} \emph{I can say a learned Telugu word for a buzz or hum, recognise its uncommon first letter, and write that letter in a source-backed order.}

Take \te{ఝ} from the already-known word \te{ఝంకారం} and write its five separately numbered movements.
\end{chapteropening}

\section[\texorpdfstring{\te{ఝంకారం} (jhaṅkāraṁ)}{ఝంకారం (jhaṅkāraṁ)}]{\te{ఝంకారం} (jhaṅkāraṁ) — a buzz, a hum}

\label{lesson:TE-C152-jhankaram}



\begin{quote}
Imagine a bee circling close enough that its sound fills the air. You already know \textbf{\te{క}}, the dot-like \textbf{\te{ం}}, and the long \textbf{\te{ా}} sign. One rare first letter will turn those familiar pieces into a name for that sound.
\end{quote}

\subsection*{You'll want to know: \te{ఝంకారం}}
\textbf{\te{ఝంకారం}} — \emph{jhaṅkāraṁ} — \textquotedblleft{}a buzz\textquotedblright{} or \textquotedblleft{}a hum\textquotedblright{}.

Say it in three pieces: \emph{jhaṅ · kā · raṁ}. The first sound is an aspirated \textbf{jh}, with a small breath after the voice. Read the word as \textbf{\te{ఝం}-\te{కా}-\te{రం}}.

This is a learned, literary word derived from Sanskrit. Older dictionaries also give it a wider sense of a loud resonant noise. Keep today's image narrow: the steady buzzing or humming sound, especially the sound of a bee.

\begin{grammarlens}[title={What you've built}]
You can hear and say the whole word before the next lesson asks your hand to isolate its uncommon first letter.
\end{grammarlens}

\subsection*{Your turn}
\begin{itemize}
\raggedright
  \item \emph{Say it:} \emph{jhaṅ · kā · raṁ}, slowly
  \item \emph{Say it:} \emph{jhaṅkāraṁ}, as one word
  \item \emph{Read it:} \textbf{\te{ఝంకారం}} — a buzz or hum
  \item \emph{Find:} familiar \textbf{\te{కా}} in the middle
\end{itemize}

\begin{tcolorbox}[breakable,colback=teal!4,colframe=teal!35!black,title={Before you move on}]
What does \textbf{\te{ఝంకారం}} mean? (\textquotedblleft{}A buzz\textquotedblright{} or \textquotedblleft{}a hum.\textquotedblright{}) Say it once more.
\end{tcolorbox}
\section[\texorpdfstring{\te{ఝ} (jha)}{ఝ (jha)}]{\te{ఝ} — the breathed first letter in \te{ఝంకారం}}

\label{lesson:TE-S173-letter-jha}



\begin{quote}
Say \textbf{\te{ఝంకారం}} once. What sound can it name? (A buzz or hum.)

Now slow the first piece: \emph{jhaṅ}. The base letter carrying that breathed \textbf{jh} is today's whole task.
\end{quote}

\subsection*{Script you'll notice: \te{ఝ}}
\textbf{\te{ఝ}} — \emph{jha}.

In \textbf{\te{ఝంకారం}}, the dot-like \textbf{\te{ం}} closes the first syllable: \textbf{\te{ఝం}}. Remove that sign and the base letter is \textbf{\te{ఝ}}.

This aspirated letter is uncommon in everyday Telugu vocabulary. That is why the book gave you \textbf{\te{ఝంకారం}} before asking you to practise the shape.

\subsection*{Writing: \te{ఝ} — five movements, five runs}
\hlblockfigure{\detokenize{figures/TE-S173-letter-jha-filmstrip.pdf}}{How \te{ఝ} is written, stroke by stroke}

Follow the filmstrip slowly. Movements 1–3 travel clockwise through the three round sections. Movement 4 curls upward through the flourish above them. Movement 5 draws the separate stem downward beneath the middle-right side.

The numbered route is one attested school-style order; Telugu handwriting can vary. Keep the five pen-down runs distinct and let all three bowls stay round.

\subsection*{Your turn}
\begin{itemize}
\raggedright
  \item \emph{Trace it:} \textbf{\te{ఝ}} once, following the five numbered movements
  \item \emph{Copy:} \textbf{\te{ఝ}} twice without tracing
  \item \emph{Build:} \textbf{\te{ఝ}} + \textbf{\te{ం}} = \textbf{\te{ఝం}}
  \item \emph{Read it:} \textbf{\te{ఝంకారం}}
\end{itemize}

\begin{tcolorbox}[breakable,colback=teal!4,colframe=teal!35!black,title={Before you move on}]
Which letter is this — \textbf{\te{ఝ}}? What word gave it to you? (\textbf{\te{ఝంకారం}}.)
\end{tcolorbox}
\section[\texorpdfstring{\te{ఝంకారం} (jhaṅkāraṁ)}{ఝంకారం (jhaṅkāraṁ)}]{\te{ఝంకారం} — meaning and shape together}

\label{lesson:TE-R152-jhankaram-recall}



\begin{quote}
Imagine the steady sound of a bee. Say the learned Telugu word for that buzz or hum. (\textbf{\te{ఝంకారం}}.)
\end{quote}

\subsection*{Your turn}
Read \textbf{\te{ఝంకారం}} once. Cover it, then write only its first base letter: \textbf{\te{ఝ}}.

Now uncover the word and check two things: the three round joined bowls and the short separate stem beneath the right side. Add \textbf{\te{ం}} only after the base letter is steady.

\begin{itemize}
\raggedright
  \item \emph{Write it:} \textbf{\te{ఝ}} once from memory
  \item \emph{Build:} \textbf{\te{ఝ}} + \textbf{\te{ం}} = \textbf{\te{ఝం}}
  \item \emph{Say it:} “a buzz or hum”
\end{itemize}

\begin{tcolorbox}[breakable,colback=teal!4,colframe=teal!35!black,title={Before you move on}]
What does \textbf{\te{ఝంకారం}} mean, and which base letter begins it? (\textquotedblleft{}A buzz\textquotedblright{} or \textquotedblleft{}a hum\textquotedblright{}; \textbf{\te{ఝ}}.)
\end{tcolorbox}
